\subsection{Simplifying Solutions to Power Equations} 
When we solve equations using the root properties, we often get solutions that have very specific forms.
\begin{Example}
Solve $(6x-12)^4=162$
\begin{lpic}{}{4}
\end{lpic}
\end{Example}
\noi It seems like we should be able to simplify this more! Can we ``cancel'' that 3?
\begin{definition}
If we have an expression in the form
$$\frac{a+b\sqrt[p]{c}}{d}$$
it is considered to be in its \term{simplest form} if the radical expression $b\sqrt[p]{c}$ is simplified, the denominator is rationalized, and $a$, $b$, and $d$ have no factors in common.
\end{definition}
To cancel any factors shared by $a$, $b$, and $d$, we treat this like a rational expression.
\begin{Example}Simplify $\dfrac{12\pm 3\sqrt[4]{2}}{6}$.
\begin{lpic}{}{2}
\fractextleft{0}{-1}{$\dfrac{12\pm 3\sqrt[4]{2}}{6}={}$};
\end{lpic}
\end{Example}
\begin{Note}We \emph{cannot} cancel a factor that appears in only one term of the numerator, just like with rational expressions.\end{Note}
\begin{Example}Solve $3(2x-6)^2+14=230$. Simplify your answer as much as possible.
\begin{lpic}{}{6}
\end{lpic}
\end{Example}